\addsectionframe{2. Anforderungen und Projektablauf}

\begin{itemframe}{Anforderungen als Leitplanken}
  \item Didaktik: nicht nur Ergebnisse, sondern Zwischenschritte sichtbar machen.
  \item Bedienung: gleiche Navigation, Tabs und Rückmeldungen in allen Themen.
  \item Technik: neue Themen sollen ohne Änderungen an Navigation und Routing ergänzt werden können.
  \item Qualität: Deutsch und Englisch, helles und dunkles Erscheinungsbild sowie automatisierte Prüfungen.
\end{itemframe}

\begin{itemframe}{Anforderungsprüfung mit Lehrenden}
  \item \textbf{Hier den tatsächlichen Nachweis ergänzen:} Gesprächspartner, Datum und Anlass.
  \item Beispielhafte Folienstruktur: Anforderung aus dem Gespräch \(\rightarrow\) Umsetzung \(\rightarrow\) beobachtetes Ergebnis.
  \item Geeignete Belege: Gesprächsnotiz, freigegebene E-Mail, Protokoll oder Feedback zu einem Prototypen.
  \item Ohne diesen Beleg sollte die Präsentation keine Kundengespräche behaupten.
\end{itemframe}

\begin{itemframe}{Projektablauf}
  \item Anforderungen und Themenumfang aus Lehrmaterial und Lastenheft ableiten.
  \item Architektur und gemeinsame UI-Bausteine vor den Fachmodulen festlegen.
  \item Fachmodule in Feature-Branches entwickeln und per Merge Request integrieren.
  \item Tests, Build und Lighthouse-Prüfung in GitLab ausführen; stabile Stände nach \texttt{main} überführen.
\end{itemframe}
